../cictikz/src/cictikz/data/tex/cictikz_preamble.tex

% Two ways to wire a two-bit resistor-string DAC, side by side. Both have
% four resistors; they differ only in which four nodes the first rank of
% switches taps.
%
% Left  — the top tap is V_REF itself, so the output runs 1/4 .. 4/4 V_REF.
% Right — the top tap is one resistor down and the bottom tap is ground, so
%         the output runs 0/4 .. 3/4 V_REF.
%
% That is the question the lecture asks under the figure, so the two halves
% are drawn by one macro and differ only in the index of the topmost tap.

\begin{circuitikz}[american, thick, transform shape, circuit ee IEC]

  ../cictikz/src/cictikz/data/tex/cictikz_lib.tex
  ../cictikz/src/cictikz/data/tex/rdac_lib.tex

  % #1 x origin, #2 index of the topmost tap used, #3 node name prefix.
  % The taps are laid down as named coordinates first: TikZ wants a literal
  % (...) after `at`, so a macro that expands to a coordinate cannot be
  % passed through another macro.
  \newcommand{\rdacHalf}[3]{%
    % Pitch 2.1: a \vresistor is \grid = 1.6 tall, and the extra half unit
    % leaves room for a switch and its control label between two taps.
    \foreach \i in {0,...,4} {
      \coordinate (#3t\i) at ({#1}, {0.5 + 2.1*\i});
    }
    \pgfmathtruncatemacro{\tapA}{#2}
    \pgfmathtruncatemacro{\tapB}{#2-1}
    \pgfmathtruncatemacro{\tapC}{#2-2}
    \pgfmathtruncatemacro{\tapD}{#2-3}
    %
    % ---- The string ----
    \draw (#3t0) \vresistor{$R$} -- ++(0,0.5)
                 \vresistor{$R$} -- ++(0,0.5)
                 \vresistor{$R$} -- ++(0,0.5)
                 \vresistor{$R$} -- ++(0,0.5);
    \draw (#3t0) -- ++(0,-0.3) \vground;
    \draw (#3t4) -- ++(-1.0,0) coordinate (#3vref);
    \rdacPort{(#3vref)}{south}{$V_{REF}$}
    %
    % ---- First rank: b_0 and its complement, twice ----
    \rdacSw{#3MA}{(#3t\tapA)}{2.0}{$b_0$}
    \rdacSw{#3MB}{(#3t\tapB)}{2.0}{$\overline{b_0}$}
    \rdacSw{#3MC}{(#3t\tapC)}{2.0}{$b_0$}
    \rdacSw{#3MD}{(#3t\tapD)}{2.0}{$\overline{b_0}$}
    %
    % Each pair shares an output, halfway between its two taps.
    \coordinate (#3pa) at ($(#3t\tapA)!0.5!(#3t\tapB)+(2.0,0)$);
    \coordinate (#3pb) at ($(#3t\tapC)!0.5!(#3t\tapD)+(2.0,0)$);
    \draw ($(#3t\tapA)+(2.0,0)$) -- ($(#3t\tapB)+(2.0,0)$);
    \draw ($(#3t\tapC)+(2.0,0)$) -- ($(#3t\tapD)+(2.0,0)$);
    \fill (#3pa) circle (0.07);
    \fill (#3pb) circle (0.07);
    %
    % ---- Second rank: b_1 picks which pair reaches the output ----
    \draw (#3pa) -- ++(0.8,0);
    \draw (#3pb) -- ++(0.8,0);
    \coordinate (#3qa) at ($(#3pa)+(0.8,0)$);
    \coordinate (#3qb) at ($(#3pb)+(0.8,0)$);
    \rdacSw{#3ME}{(#3qa)}{2.0}{$b_1$}
    \rdacSw{#3MF}{(#3qb)}{2.0}{$\overline{b_1}$}
    \draw ($(#3qa)+(2.0,0)$) -- ($(#3qb)+(2.0,0)$);
    \coordinate (#3vo) at ($($(#3qa)!0.5!(#3qb)$)+(2.0,0)$);
    \fill (#3vo) circle (0.07);
    \draw (#3vo) -- ++(0.8,0) coordinate (#3voport);
    \rdacPort{(#3voport)}{south west}{$V_{OUT}$}
  }

  \rdacHalf{0}{4}{L}
  \rdacHalf{10.5}{3}{R}

\end{circuitikz}
\end{document}
